\usetikzlibrary{shapes.geometric,shapes.misc}

\begin{figure}[H]
\centering
\begin{tikzpicture}[
x=.35cm,
y=.35cm,
edge/.style={draw=black!55, line width=.25pt},
circlev/.style={circle, fill=black, draw=black, inner sep=1.4pt},
starv/.style={star, star points=5, star point ratio=2.1, fill=black, draw=black, inner sep=0pt, minimum size=6pt},
crossv/.style={cross out, draw=black, line width=.6pt, inner sep=0pt, minimum size=5pt},
crossv /.style={cross out, draw=black, line width=.6pt, inner sep=0pt, minimum size=5pt},
squarev/.style={rectangle, fill=black, draw=black, inner sep=1.4pt}
]

\def\nodedata{
n1/1/25/circlev,
n2/5/23/circlev,
n3/6/30/circlev,
n4/7/21/circlev,
n5/8/28/circlev,
n7/10/26/circlev,
n8/12/24/circlev,
n10/9/17/starv,
n11/13/15/starv,
n12/14/22/starv,
n13/15/13/starv,
n14/16/20/starv,
n16/18/18/starv,
n17/20/16/starv,
n19/17/9/crossv,
n20/21/7/crossv,
n21/22/14/crossv,
n22/23/5/crossv,
n24/25/1/squarev,
n25/26/10/crossv,
n26/28/8/crossv,
n27/30/6/squarev,
n28/24/12/crossv
}

\foreach[count=\i] \name/\xx/\yy/\style in \nodedata {
    \coordinate (\name) at (\xx,\yy);
    \expandafter\xdef\csname X\i\endcsname{\xx}
    \expandafter\xdef\csname Y\i\endcsname{\yy}
    \expandafter\xdef\csname Name\i\endcsname{\name}
    \xdef\NNodes{\i}
}

\pgfmathtruncatemacro{\LastButOne}{\NNodes-1}
\foreach \i in {1,...,\LastButOne} {
    \pgfmathtruncatemacro{\jstart}{\i+1}
    \foreach \j in {\jstart,...,\NNodes} {
        \edef\Xi{\csname X\i\endcsname}
        \edef\Yi{\csname Y\i\endcsname}
        \edef\Xj{\csname X\j\endcsname}
        \edef\Yj{\csname Y\j\endcsname}
        \edef\NameI{\csname Name\i\endcsname}
        \edef\NameJ{\csname Name\j\endcsname}
        \pgfmathtruncatemacro{\cmp}{((\Xj>=\Xi)&&(\Yj>=\Yi)) || ((\Xi>=\Xj)&&(\Yi>=\Yj))}
        \ifnum\cmp=1
            \draw[edge] (\NameI) -- (\NameJ);
        \fi
    }
}

\foreach \name/\xx/\yy/\style in \nodedata {
    \node[\style] at (\name) {};
}

\end{tikzpicture}
\caption{Three tiled copies of the diagonal graph gadget. The first copy excluding its right boundary is shown in $\bullet$. The second copy excluding its right boundary is shown in $\filledstar$. The third copy excluding its right boundary is shown in $\times$. The third copy's right boundary is shown in $\filledsquare$.}
\end{figure}